% !TEX root = ../../ctfp-print.tex

\lettrine[lhang=0.17]{一}{个范畴是小}的，如果它的对象构成一个集合。但我们知道存在比集合更大的事物。著名的是，在标准集合论（策梅洛-弗兰克尔理论，可选地附加选择公理）中无法构造出包含所有集合的集合。因此所有集合的范畴必然是大的。数学上有一些技巧如格罗滕迪克宇宙（Grothendieck universes），可以用来定义超越普通集合的集合论系统。这些技巧使我们能够讨论大范畴。

一个范畴称为\emph{局部小范畴}，如果任意两个对象间的态射构成一个集合。当它们不构成集合时，我们需要重新思考若干定义。特别地，当我们无法从集合中选取态射时，组合态射意味着什么？解决方案是通过替换操作进行自举：将作为$\Set$中对象的Hom集替换为某个其他范畴$\cat{V}$中的\emph{对象}。区别在于一般情况下对象没有元素，因此我们不能再谈论个别态射。必须用能够对同态对象整体执行的操作来重新定义\emph{丰富化}范畴的所有性质。为此提供同态对象的范畴必须具有额外结构——它必须是一个单子范畴。若我们将此单子范畴记作$\cat{V}$，就可以讨论在$\cat{V}$上的范畴$\cat{C}$。

除了大小原因外，我们可能还希望将Hom集推广到比普通集合更具结构的对象。例如，传统范畴不具备对象间距离的概念。两个对象要么由态射连接，要么没有连接。所有与给定对象相连的对象都是它的邻居。这与现实不同；在范畴中，朋友的朋友的朋友与密友同样亲近。在适当丰富化的范畴中，我们可以定义对象间的距离。

还有一个非常实用的理由让我们需要获得丰富化范畴的经验：因为一个非常有用的范畴论在线知识库\urlref{https://ncatlab.org/}{nLab}主要就是用丰富化范畴的语言书写的。

\section{为何需要单子范畴？}

构造丰富化范畴时，我们需要牢记：当把单子范畴替换为$\Set$并将同态对象替换为Hom集时，应当能够恢复常规定义。实现这一目标的最佳方式是从常规定义出发，以"无点"方式（即不涉及集合元素名称的方式）持续重构它们。

让我们从组合律的定义开始。通常情况下，它取一对态射：一个来自$\cat{C}(b, c)$，另一个来自$\cat{C}(a, b)$，并将其映射到$\cat{C}(a, c)$中的态射。换句话说，这是一个映射：
$$\cat{C}(b, c)\times{}\cat{C}(a, b) \to \cat{C}(a, c)$$
这是集合之间的函数——其中一个集合是两个Hom集的笛卡尔积。通过用更一般的结构替代笛卡尔积，这个公式可以轻松推广。范畴积是可行的，但我们甚至可以用完全一般的张量积。

接下来考虑恒等态射。与其从Hom集中选取单独元素，不如用从单元素集合$\cat{1}$出发的函数来定义它们：
$$j_a \Colon \cat{1} \to \cat{C}(a, a)$$
同样，我们可以用终对象替代单元素集合，但更进一步的做法是用张量积的单位元$i$来替代。

如您所见，从某个单子范畴$\cat{V}$中取出的对象很适合作为Hom集的替代品。

\section{单范畴}

我们之前讨论过单范畴，但重新陈述其定义是有意义的。单范畴定义了一个作为双函子的张量积：
$$\otimes \Colon \cat{V}\times{}\cat{V} \to \cat{V}$$
我们希望张量积具有结合律，但只需满足到自然同构的程度即可。这种同构称为结合子（associator），其分量为：
$$\alpha_{a b c} \Colon (a \otimes b) \otimes c \to a \otimes (b \otimes c)$$
它必须在所有三个变元中自然成立。

单范畴还必须定义一个特殊的单位对象 $i$，作为张量积的单位元；同样地，该性质只需满足到自然同构。这两个同构分别称为左单位子和右单位子（left/right unitor），其分量为：
\begin{align*}
  \lambda_a & \Colon i \otimes a \to a \\
  \rho_a    & \Colon a \otimes i \to a
\end{align*}
结合子与单位子必须满足如下相容条件：

\begin{figure}[H]
  \centering
  \begin{tikzcd}[row sep=large]
    ((a \otimes b) \otimes c) \otimes d
    \arrow[d, "\alpha_{(a \otimes b)cd}"]
    \arrow[rr, "\alpha_{abc} \otimes \id_d"]
    & & (a \otimes (b \otimes c)) \otimes d
    \arrow[d, "\alpha_{a(b \otimes c)d}"] \\
    (a \otimes b) \otimes (c \otimes d)
    \arrow[rd, "\alpha_{ab(c \otimes d)}"]
    & & a \otimes ((b \otimes c) \otimes d)
    \arrow[ld, "\id_a \otimes \alpha_{bcd}"] \\
    & a \otimes (b \otimes (c \otimes d))
  \end{tikzcd}
\end{figure}

\begin{figure}[H]
  \centering
  \begin{tikzcd}[row sep=large]
    (a \otimes i) \otimes b
    \arrow[dr, "\rho_{a} \otimes \id_b"']
    \arrow[rr, "\alpha_{aib}"]
    & & a \otimes (i \otimes b)
    \arrow[dl, "\id_a \otimes \lambda_b"] \\
    & a \otimes b
  \end{tikzcd}
\end{figure}

\noindent
若存在分量为
$$\gamma_{a b} \Colon a \otimes b \to b \otimes a$$
的自然同构，并且满足「平方等于恒等」条件：
$$\gamma_{b a} \circ \gamma_{a b} = \idarrow[a \otimes b]$$
同时与单结构相容，则该单范畴称为\newterm{对称单范畴（symmetric monoidal category）}。

单范畴的一个有趣特性是，你可能可以将内态同态（函数对象）定义为张量积的右伴随。你或许记得，函数对象（即指数对象）的标准定义是通过范畴积的右伴随给出的。对于任意对象对都存在这种对象的范畴被称为笛卡尔闭范畴。以下伴随关系定义了单范畴中的内态同态：
$$\cat{V}(a \otimes b, c) \sim \cat{V}(a, [b, c])$$
根据
\urlref{http://www.tac.mta.ca/tac/reprints/articles/10/tr10.pdf}{G. M.
  Kelly}的记号约定，我采用${[}b, c{]}$表示内态同态。该伴随的余单位是一个自然变换，其分量称为求值态射：
$$\varepsilon_{a b} \Colon ([a, b] \otimes a) \to b$$
请注意，若张量积不对称，我们可以使用如下伴随定义另一个内态同态${[}{[}a, c{]}{]}$：
$$\cat{V}(a \otimes b, c) \sim \cat{V}(b, [[a, c]])$$
同时定义这两种内态同态的单范畴称为\newterm{双闭（biclosed）}。非双闭范畴的一个典型例子是集合范畴$\Set$上的自函子范畴，其张量积由函子复合构成——这正是我们用来定义单子（monad）的那个范畴。

\section{丰富范畴}

一个在单子范畴 $\cat{V}$ 上丰富的范畴 $\cat{C}$ 将同态集替换为同态对象。对于 $\cat{C}$ 中的每对对象 $a$ 和 $b$，我们关联一个 $\cat{V}$ 中的对象 $\cat{C}(a, b)$。我们使用与同态集相同的记号表示同态对象（hom-object），但需明确它们并不包含态射。另一方面，$\cat{V}$ 是具有同态集和态射的常规（非丰富化）范畴。因此我们并未彻底摆脱集合——只是将它们隐藏到了底层结构之下。

由于无法讨论 $\cat{C}$ 中的单个态射，态射的合成被 $\cat{V}$ 中的态射族取代：
$$\circ \Colon \cat{C}(b, c) \otimes \cat{C}(a, b) \to \cat{C}(a, c)$$

\begin{figure}[H]
  \centering
  \includegraphics[width=0.45\textwidth]{images/composition.jpg}
\end{figure}

\noindent
类似地，单位态射被替换为 $\cat{V}$ 中的态射族：
$$j_a \Colon i \to \cat{C}(a, a)$$
其中 $i$ 是 $\cat{V}$ 中的张量单位。

\begin{figure}[H]
  \centering
  \includegraphics[width=0.4\textwidth]{images/id.jpg}
\end{figure}

\noindent
结合律通过 $\cat{V}$ 中的结合约束（associator）定义：

\begin{figure}[H]
  \centering
  \begin{tikzcd}[column sep=large]
    (\cat{C}(c,d) \otimes \cat{C}(b,c)) \otimes \cat{C}(a,b)
    \arrow[r, "\circ\otimes\id"]
    \arrow[dd, "\alpha"]
    & \cat{C}(b,d) \otimes \cat{C}(a,b)
    \arrow[dr, "\circ"] \\
    & & \cat{C}(a,d) \\
    \cat{C}(c,d) \otimes (\cat{C}(b,c) \otimes \cat{C}(a,b))
    \arrow[r, "\id\otimes\circ"]
    & \cat{C}(c,d) \otimes \cat{C}(a,c)
    \arrow[ur, "\circ"]
  \end{tikzcd}
\end{figure}

\noindent
单位律同样通过单位约束（unitor）表达：

\begin{figure}[H]
  \centering
  \begin{subfigure}
    \centering
    \begin{tikzcd}[row sep=large]
      \cat{C}(a,b) \otimes i
      \arrow[rr, "\id \otimes j_a"]
      \arrow[dr, "\rho"]
      & & \cat{C}(a,b) \otimes \cat{C}(a,a)
      \arrow[dl, "\circ"] \\
      & \cat{C}(a,b)
    \end{tikzcd}
  \end{subfigure}
  \hspace{1cm}
  \begin{subfigure}
    \centering
    \begin{tikzcd}[row sep=large]
      i \otimes \cat{C}(a,b)
      \arrow[rr, "j_b \otimes \id"]
      \arrow[dr, "\lambda"]
      & & \cat{C}(b,b) \otimes \cat{C}(a,b)
      \arrow[dl, "\circ"] \\
      & \cat{C}(a,b)
    \end{tikzcd}
  \end{subfigure}
\end{figure}

\section{预序}

预序定义为一个薄范畴，即其中每个态射集要么为空要么为单元素集合。我们将非空的态射集 $\cat{C}(a, b)$ 解释为 $a$ 小于等于 $b$ 的证明。这种范畴可以视为在仅包含两个对象 $0$ 和 $1$（有时称为 $\mathit{False}$ 和 $\mathit{True}$）的极简幺半范畴上的丰富化范畴。除了必要的恒等态射外，该范畴还包含唯一的从 $0$ 到 $1$ 的态射 $0 \to 1$。通过将张量积定义为 $0$ 和 $1$ 的简单算术运算（即唯一的非零乘积为 $1 \otimes 1$），可建立简单的幺半结构。该范畴的单位对象是 $1$，且这是一个严格幺半范畴，即结合子和单位子均为恒等态射。

由于在预序中态射集要么为空要么为单元素集，我们可以轻易地用这个微小范畴中的同调对象替代它。富化的预序 $\cat{C}$ 对任意对象对 $a,b$ 都具有同调对象 $\cat{C}(a, b)$。若 $a \leqslant b$，该对象为 $1$；否则为 $0$。

观察组合性质：任意两个对象的张量积为 $0$，除非两者均为 $1$ 时结果为 $1$。当结果为 $0$ 时，组合态射有两种选择：$\idarrow[0]$ 或 $0 \to 1$；而结果为 $1$ 时只能选择 $\idarrow[1]$。将其还原为关系语言，这表示若 $a \leqslant b$ 且 $b \leqslant c$，则必有 $a \leqslant c$，这正是预序所需的传递律。

关于恒等性：它是从 $1$ 到 $\cat{C}(a, a)$ 的态射。从 $1$ 出发仅有恒等态射 $\idarrow[1]$，因此 $\cat{C}(a, a)$ 必为 $1$，这意味着 $a \leqslant a$，即预序的自反律。由此可见，若将预序实现为丰富化范畴，则传递性和自反性会自动满足。

\section{度量空间}

一个有趣的例子归功于
\urlref{http://www.tac.mta.ca/tac/reprints/articles/1/tr1.pdf}{William
  Lawvere}。他注意到度量空间可以通过充实范畴(enriched categories)来定义。度量空间定义了任意两个对象之间的距离，这种距离是非负实数。我们通常会将无穷大作为可能的取值包含进来，当距离为无穷大时，表示无法从起点对象到达目标对象。

距离函数需要满足若干明显性质：其中之一是对象到自身的距离必须为零；另一个是三角不等式(triangle inequality)，即直接距离不大于经过中间点的距离之和。我们不要求距离具有对称性，这初看可能显得奇怪，正如Lawvere所解释的那样，可以想象在一个方向上是上坡（需要更大代价），而在另一方向则是下坡。当然，如有必要，后续可以额外施加对称性约束。

如何将度量空间转化为范畴论语言？我们需要构造一个范畴，其hom-对象就是距离。请注意，这里的距离不是态射而是hom-对象。如何让hom-对象成为数值？这需要构造一个单子范畴(monoidal category) $\cat{V}$，其中这些数值作为对象存在。非负实数（加上无穷大）构成全序关系，因此可以视为薄范畴(thin category)。两个数值$x$和$y$之间存在态射当且仅当$x \geqslant y$（注意：这与预序关系传统定义的方向相反）。单子结构由加法给出，零元作为单位对象。换句话说，两个数值的张量积就是它们的和。

度量空间是相对于此类单子范畴的充实范畴。从对象$a$到$b$的hom-对象$\cat{C}(a, b)$是一个非负（可能无穷）数值，我们称之为从$a$到$b$的距离。让我们考察此类范畴中的恒等和合成条件。

根据定义，从张量单位（即数值0）到hom-对象$\cat{C}(a, a)$的态射对应关系：
$$0 \geqslant \cat{C}(a, a)$$
由于$\cat{C}(a, a)$是非负数值，此条件表明从$a$到自身的距离恒为零。
验证通过！

现在讨论合成条件。我们从两个相邻hom-对象的张量积开始：$\cat{C}(b, c) \otimes \cat{C}(a, b)$。根据定义张量积为两距离之和。合成是$\cat{V}$中从该乘积到$\cat{C}(a, c)$的态射。$\cat{V}$中的态射定义为大于等于关系。换句话说，从$a$到$b$再从$b$到$c$的距离之和大于等于从$a$到$c$的直接距离。这正是标准的三角不等式。验证通过！

通过将度量空间重新表述为充实范畴，我们自然获得了三角不等式和零自反距离这两个性质。

\section{强化函子}

函子的定义涉及态射的映射。在强化语境中，我们不再具有单个态射的概念，
因此必须整体处理态射对象(hom-objects)。态射对象是对称单子范畴$\cat{V}$中的对象，
我们可以利用其间的态射进行操作。因此，当两个范畴均以相同的对称单子范畴$\cat{V}$为强化基底时，
我们便能定义它们之间的强化函子。此时可用$\cat{V}$中的态射来映射两个强化范畴间的态射对象。

在两个范畴$\cat{C}$与$\cat{D}$之间的\newterm{强化函子}(enriched functor) $F$，
除了将对象映射到对象外，还需为$\cat{C}$中的每对对象指定$\cat{V}$中的一个态射：
$$F_{a b} \Colon \cat{C}(a, b) \to \cat{D}(F a, F b)$$
函子本质是保持结构的映射。对于普通函子而言，这意味着保持复合运算和恒等态射。
在强化语境中，保持复合结构意味着下图交换：

\begin{figure}[H]
  \centering
  \begin{tikzcd}[column sep=large, row sep=large]
    \cat{C}(b,c) \otimes \cat{C}(a,b)
    \arrow[r, "\circ"]
    \arrow[d, "F_{bc} \otimes F_{ab}"]
    & \cat{C}(a,c)
    \arrow[d, "F_{ac}"] \\
    \cat{D}(F b, F c) \otimes \cat{D}(F a, F b)
    \arrow[r,  "\circ"]
    & \cat{D}(F a, F c)
  \end{tikzcd}
\end{figure}

\noindent
保持恒等态射的条件被替换为保持$\cat{V}$中"选择"恒等元的态射：

\begin{figure}[H]
  \centering
  \begin{tikzcd}[row sep=large]
    & i \arrow[dl, "j_a"'] \arrow[dr, "j_{F a}"] & \\
    \cat{C}(a,a)
    \arrow[rr, "F_{aa}"]
    & & \cat{D}(F a, F a)
  \end{tikzcd}
\end{figure}

\section{自丰富化}

一个闭对称单子范畴可以通过用内部Hom对象替代态射集来实现自丰富化（见上文定义）。为了实现这一点，我们需要为内部Hom对象定义合成律。换句话说，我们需要实现具有以下形式的态射：
$$[b, c] \otimes [a, b] \to [a, c]$$
这与普通编程任务差别不大，只不过在范畴论中我们通常采用无点式实现。我们首先指定该元素所属的集合，即如下Hom集：
$$\cat{V}([b, c] \otimes [a, b], [a, c])$$
这个Hom集同构于：
$$\cat{V}(([b, c] \otimes [a, b]) \otimes a, c)$$
这里刚使用了定义内部Hom $[a, c]$ 的伴随关系。如果我们能在新的Hom集中构建某个态射，该伴随关系就会指引我们找到原始Hom集中的对应态射，从而作为所需的合成律。我们通过组合若干现有态射来构造这个态射。首先，我们可以使用结合子 $\alpha_{[b, c]\ [a, b]\ a}$ 对左侧表达式重新结合：
$$([b, c] \otimes [a, b]) \otimes a \to [b, c] \otimes ([a, b] \otimes a)$$
接着应用伴随关系的余单位 $\varepsilon_{a b}$：
$$[b, c] \otimes ([a, b] \otimes a) \to [b, c] \otimes b$$
再次应用余单位 $\varepsilon_{b c}$ 即可得到 $c$。于是我们构造出如下态射：
$$\varepsilon_{b c}\ .\ (\idarrow[{[b, c]}] \otimes \varepsilon_{a b})\ .\ \alpha_{[b, c] [a, b] a}$$
它属于Hom集：
$$\cat{V}(([b, c] \otimes [a, b]) \otimes a, c)$$
通过伴随关系就得到了我们所求的合成律。

同样地，恒等态射：
$$j_a \Colon i \to [a, a]$$
属于如下Hom集：
$$\cat{V}(i, [a, a])$$
通过伴随关系该Hom集同构于：
$$\cat{V}(i \otimes a, a)$$
我们知道这个Hom集包含左单位元 $\lambda_a$。我们可以定义 $j_a$ 为该元素通过伴随关系对应的像。

自丰富化的一个实际例子是范畴 $\Set$，它是编程语言类型系统的原型。此前我们已知 $\Set$ 关于笛卡尔积构成闭单子范畴。在 $\Set$ 中，任意两个集合间的Hom集本身也是集合，因此属于 $\Set$ 的对象。我们知道它同构于指数集合，故外部Hom与内部Hom在此等价。现在我们还知道，通过自丰富化，我们可以用指数集合作为Hom对象，并以指数对象的笛卡尔积来表达合成运算。

\section{与$\cat{2}$-范畴的关系}

我在$\Cat$（小范畴的范畴）的语境下讨论过$\cat{2}$-范畴。范畴间的态射是函子，但还存在额外的结构：函子之间的自然变换。在$\cat{2}$-范畴中，对象常被称为零维胞腔(zero-cells)；态射被称为一维胞腔($1$-cells)；态射间的态射被称为二维胞腔($2$-cells)。在$\Cat$中，零维胞腔是范畴，一维胞腔是函子，二维胞腔是自然变换。

但需要注意的是，两个范畴间的函子本身也构成一个范畴；因此在$\Cat$中，我们实际上拥有一个\emph{同态范畴}(hom-category)而非同态集(hom-set)。事实上，正如$\Set$可以被视为在$\Set$上充实的范畴，$\Cat$也可以被视为在$\Cat$上充实的范畴。更一般地，正如每个范畴都可以被视为在$\Set$上充实的范畴，每个$\cat{2}$-范畴都可以被视为在$\Cat$上充实的范畴。